\section*{الفصل \ref{ch:b2:aromatic-substitution} --- العطرية والاستبدال العطري المحب للإلكترونات}

\begin{solution}{exo:b2:aromatic-substitution:1}
\ce{HNO3 + 2H2SO4 -> NO2+ + H3O+ + 2HSO4-}: يُبرتَن حمض النتريك، ثم يفقد الماء.
\end{solution}

\begin{solution}{exo:b2:aromatic-substitution:2}
يعطي التولوين 2-بروموتولوين و4-بروموتولوين: فالميثيل يوجّه إلى الموضعين أورثو 
وبارا. ويعطي النتروبنزين 3-برومونتروبنزين، ببطء أكبر: فمجموعة النترو
توجّه إلى ميتا.
\end{solution}

\begin{solution}{exo:b2:aromatic-substitution:3}
\ce{-CH3}: منشطة، أورثو/بارا. \ce{-OH}: منشطة بقوة، أورثو/بارا. \ce{-Cl}: مثبطة، أورثو/بارا.
\ce{-NO2}: مثبطة بقوة، ميتا. \ce{-COCH3}: مثبطة، ميتا. \ce{-NHCOCH3}:
منشطة (باعتدال)، أورثو/بارا.
\end{solution}

\begin{solution}{exo:b2:aromatic-substitution:4}
الأسيتوفينون (1-فينيل إيثان-1-ون)، \ce{C6H5COCH3}، عبر \omterm{def:b2:aromatic-substitution:friedel-crafts}{أيون الأسيليوم} \ce{CH3CO+}.
\end{solution}

\begin{solution}{exo:b2:aromatic-substitution:5}
يعاد ترتيب الكاتيون الكربوني الأولي (أو معقده) بانتقال هيدريد إلى كاتيون
الإيزوبروبيل الثانوي، الذي يؤلكل الحلقة. ولتحضير بروبيل بنزين: \omterm{def:b2:aromatic-substitution:friedel-crafts}{أسيلة فريدل--كرافتس} بكلوريد
البروبانويل (بلا إعادة ترتيب)، ثم اختزال \ce{C=O} الكيتون إلى \ce{CH2}.
\end{solution}

\begin{solution}{exo:b2:aromatic-substitution:6}
مجموعة \ce{-OH} مانحة ميزوميرية قوية: فالحلقة منشطة إلى حد يكفي معه البروم الجزيئي
محبًا للإلكترونات، فتُستبدل المواضع الثلاثة أورثو وبارا كلها.
\end{solution}

\begin{solution}{exo:b2:aromatic-substitution:7}
3-برومونتروبنزين: نترِت، ثم برومن (النترو يوجّه إلى ميتا). 4-برومونتروبنزين:
برومن، ثم نترِت (البروم يوجّه إلى أورثو/بارا)، وافصل المتماكب بارا عن
المتماكب أورثو.
\end{solution}

\begin{solution}{exo:b2:aromatic-substitution:8}
خليط النترتة يؤكسد الأنيلين الغني جدًا بالإلكترونات، ويبرتن نيتروجينه:
و\ce{-NH3+} \omterm{def:b2:aromatic-substitution:directing}{مجموعة مثبطة} موجّهة إلى ميتا. أما في الأسيتانيليد فيتقاسم الزوج الحر للنيتروجين
مع الكربونيل: فهو أقل قاعدية وأقل تنشيطًا، لا يُبرتَن ولا يتأكسد،
ويظل يوجّه إلى أورثو/بارا، وإلى بارا أساسًا لأسباب فراغية.
\end{solution}

\begin{solution}{exo:b2:aromatic-substitution:9}
سلفِن التولوين (في الموضع بارا أساسًا، إذ تتجنب المجموعة الضخمة أورثو)؛ ثم برومن: يدخل البروم
في الموضع أورثو بالنسبة إلى الميثيل (فالموضع بارا محجوب، والمجموعتان كلتاهما تفضّلان ذلك الموضع)؛
ثم أزل المجموعة \ce{SO3H} بالتسخين في حمض مخفف.
\end{solution}

\begin{solution}{exo:b2:aromatic-substitution:10}
المجموعتان كلتاهما توجّهان إلى أورثو/بارا؛ وموضع بارا لكل منهما تشغله الأخرى. فيفوز الميثوكسي،
المنشط الأقوى: نترتة في الموضع أورثو بالنسبة إلى \ce{-OCH3} (الموضع 2)، معطية
4-ميثيل-2-نتروأنيسول.
\end{solution}

\begin{solution}{exo:b2:aromatic-substitution:11}
كل مجموعة نترو تثبط الحلقة بقوة: فتحتاج النترتة الثانية إلى ظروف أقوى
من الأولى، والثالثة (على حلقة تحمل مجموعتي نترو والميثيل المنشط الضعيف
وحده) إلى ظروف أقوى بكثير.
\end{solution}

\begin{solution}{exo:b2:aromatic-substitution:12}
حمض 4-نتروبنزويك: نترِت التولوين، وافصل المتماكب بارا، وأكسد الميثيل إلى
\ce{-COOH} (البرمنغنات الساخنة). حمض 3-نتروبنزويك: أكسد التولوين إلى حمض البنزويك أولًا، ثم
نترِت (\ce{-COOH} يوجّه إلى ميتا).
\end{solution}

\begin{solution}{pb:b2:aromatic-substitution:1}
\textbf{1.} \ce{HNO3 + 2H2SO4 -> NO2+ + H3O+ + 2HSO4-}.
\textbf{2.} \ce{O=N+=O}: خطي، متساوي الإلكترونات مع \ce{CO2}.
\textbf{3.} يرتبط النيتروجين بكربون الحلقة الواقع في الموضع بارا بالنسبة إلى \ce{CH3}، فيصير رباعي الأوجه
(حاملًا H و\ce{NO2}).
\textbf{4.} الشحنة الموجبة على الكربونين أورثو بالنسبة إلى الكربون المهاجَم وعلى
الكربون بارا بالنسبة إليه، الحامل للميثيل.
\textbf{5.} الأولى، تكوّن \omterm{def:b2:aromatic-substitution:seAr}{وسيط ويلاند}.
\textbf{6.} قاعدة من الوسط: \ce{HSO4-} أو الماء.
\textbf{7.} موضعان أورثو، وموضعان ميتا، وموضع بارا واحد.
\textbf{8.} 2 : 2 : 1، أي 40~\% أورثو، و40~\% ميتا، و20~\% بارا.
\textbf{9.} في الهجوم ميتا، لا تقع الشحنة الموجبة أبدًا على الكربون الحامل للميثيل؛ فلا
تستفيد أي بنية من الميثيل.
\textbf{10.} في إحدى بنى الرنين تقع الشحنة على الكربون الحامل للمجموعة \ce{CH3}: كاتيون كربوني
ثالثي، يثبّته التأثير التحريضي للميثيل وفرط اقترانه.
\textbf{11.} أورثو وبارا مفضّلان (96~\% مقابل 60~\% في الهجوم العشوائي)؛ وميتا يكاد يغيب.
\textbf{12.} لكل موضع: أورثو $58/2 = 29~\%$، وبارا 38~\%: فبارا هو الموضع الأكثر تفاعلية،
لأن الميثيل يعيق الموضعين أورثو قليلًا.
\textbf{13.} أسرع: فالميثيل ينشط الحلقة.
\textbf{14.} 4-نتروتولوين (ينصهر عند \qty{52}{\degreeCelsius}).
\textbf{15.} برّد الخليط الخام: فيتبلور المتماكب بارا ويُرشَّح؛ ويحتفظ
السائل بالمتماكبين أورثو وميتا.
\textbf{16.} الجزيء المتناظر يتراص في البلورة على نحو أفضل: تآثرات أكثر لكل وحدة حجم،
ودرجة انصهار أعلى.
\textbf{17.} \omterm{def:b2:liquid-vapour-diagrams:fractional-distillation}{بالتقطير التجزيئي} تحت ضغط مخفَّض، أو بالكروماتوغرافيا؛ أو بمزيد من
التبريد (المتماكب ميتا ينصهر عند \qty{16}{\degreeCelsius}).
\textbf{18.} درجات غليان المتماكبات متقاربة (الصيغة نفسها، وقطبية متشابهة).
\textbf{19.} \ce{C7H8}: \qty{92.14}{g/mol}؛ \ce{C7H7NO2}: \qty{137.14}{g/mol}.
\textbf{20.} $100/92.14 = \qty{1.085}{mol}$ من التولوين؛ $0.90 \times 1.085 = \qty{0.977}{mol}$ من
النتروتولوينات.
\textbf{21.} المجموع $0.977 \times 137.14 = \qty{134.0}{g}$: أورثو \qty{77.7}{g}، وميتا
\qty{5.4}{g}، وبارا \qty{50.9}{g}.
\textbf{22.} نترتة ثانية (ثنائي نتروتولوينات)، وأكسدة مجموعة الميثيل.
\textbf{23.} أكسد الميثيل إلى \ce{-COOH} بالبرمنغنات القلوية الساخنة، ثم حمّض.
\textbf{24.} تعطي التجربة $\boldsymbol{\approx \qty{51}{g}}$ من 4-نتروتولوين.
\end{solution}
